\originalchapter{夸克、强子与深度非弹性散射}\label{ch:qcd}
\originalsection{颜色与强子产生}
电子正电子可经光子产生夸克反夸克. 在最低阶且远离共振、阈值, 截面比
$R=\sigma(e^+e^-\to\mathrm{hadrons})/\sigma(e^+e^-\to\mu^+\mu^-)=N_c\sum_f Q_f^2$, 其中开放味 $f$ 的电荷以 $e$ 为单位. 三种颜色使每味贡献乘 $N_c=3$. $u,d,s$ 开放时 $R=2$, 再开放 charm 后为 $10/3$, bottom 后为 $11/3$. 这些是最低阶平台值, 实际有 QCD 修正和共振峰.

QCD 规范群为 SU(3), 胶子耦合 $g_s\bar q\gamma^\mu T^aqG_\mu^a$. 八个生成元满足 $\Tr(T^aT^b)=\delta^{ab}/2$. 胶子自身携颜色, 因而有三胶子与四胶子顶点; 这是同 QED 光子不自耦合的重要差别. 颜色指数不是可以直接在探测器里看到的三种颜色光.

\originalsection{弹性形状因子与非弹性变量}
若电子弹性散射于质子, 质子保持完整, 电磁顶点以 Dirac 与 Pauli 形状因子 $F_1,F_2$ 描述复合结构, 常组合为 Sachs 电、磁形状因子. Fourier 意义的电荷密度解释只在适当非相对论或静止近似中直接成立, 不能对任意 relativistic 分布硬套.

8.701第5.2讲的定量入口是守恒质子电磁流。取$q=P'-P$、$Q^2=-q^2>0$，两个独立的在壳结构定义为
\begin{equation}
\langle P'|J^\mu|P\rangle
=\bar u(P')\left[F_1(Q^2)\gamma^\mu+
\frac{\ii\sigma^{\mu\nu}q_\nu}{2M_p}F_2(Q^2)\right]u(P),
\quad \sigma^{\mu\nu}=\frac{\ii}{2}[\gamma^\mu,\gamma^\nu].
\end{equation}
质子不是基本点粒子，因此只用一个常数顶点会漏掉自旋磁矩和动量转移依赖。设$\tau=Q^2/(4M_p^2)$，组合
$G_E=F_1-\tau F_2$、$G_M=F_1+F_2$；电荷归一化为$F_1(0)=G_E(0)=1$，$F_2(0)=\kappa_p$，$G_M(0)=1+\kappa_p=\mu_p/\mu_N$。静态意义下$G_E$主导电响应、$G_M$主导磁响应；它们不是两个任意可互换的拟合字母。

在单光子交换、非极化弹性$ep$、已处理靶反冲的实验室运动学下，Rosenbluth式可写
\begin{equation}
\frac{\dd\sigma}{\dd\Omega}
=\left(\frac{\dd\sigma}{\dd\Omega}\right)_{\!\rm Mott}
\frac{E'}E\left[
\frac{G_E^2+\tau G_M^2}{1+\tau}
+2\tau G_M^2\tan^2\frac\theta2\right],\quad
\epsilon=\left[1+2(1+\tau)\tan^2\frac\theta2\right]^{-1}.
\end{equation}
此处Mott因子采用上一章的\\emph{固定重靶}定义，因而把反冲通量$E'/E$显式写出；若文献把它并入Mott因子，不可再乘一次。将观测截面除以前置因子并乘$\epsilon(1+\tau)$，得到
\begin{equation}
\sigma_R\equiv\epsilon(1+\tau)
\frac{\dd\sigma/\dd\Omega}
{(\dd\sigma/\dd\Omega)_{\rm Mott}(E'/E)}
=\epsilon G_E^2+\tau G_M^2.
\end{equation}
因此固定$Q^2$而改变束能与角度（即$\epsilon$）时，斜率给$G_E^2$，截距给$\tau G_M^2$；只有一组角度的截面一般不能独立分离二者。低$Q^2$下用归一化电形状因子的斜率定义
\begin{equation}
G_E(Q^2)=1-\frac{Q^2}{6}\langle r_E^2\rangle+O(Q^4),
\qquad \langle r_E^2\rangle=-6\,\left.\frac{\dd G_E}{\dd Q^2}\right|_{0}
\quad(\hbar=c=1).
\end{equation}
若$Q^2$以GeV$^2$输入，右端须乘$(\hbar c)^2$才成为fm$^2$。磁半径则对$G_M(Q^2)/G_M(0)$取同样的负斜率。实验上须外推到$Q^2=0$并处理辐射、两光子交换等修正；不能把一个有限$Q^2$点直接当成半径。

深度非弹性散射 $e(k)+p(P)\to e(k')+X$ 中 $q=k-k'$, 定义
\begin{equation}
 Q^2=-q^2>0,\qquad x=\frac{Q^2}{2P\cdot q},\qquad y=\frac{P\cdot q}{P\cdot k},\qquad W^2=(P+q)^2.
\end{equation}
由此 $W^2=M_p^2+Q^2(1/x-1)$, 高能质量可忽略时 $Q^2=xys$. 弹性对应 $x=1$, 非弹性通常 $0<x<1$. $Q^2$ 指空间分辨尺度, 较大值分辨较短距离; $x$ 在 parton 模型中近似是被击中 parton 的纵向动量份额.

\originalsection{Parton 分布与标度}
若硬碰撞发生于近似自由 spin-$1/2$ parton, 结构函数满足最低阶
\begin{equation}
 F_2(x)=x\sum_fQ_f^2[q_f(x)+\bar q_f(x)],\qquad F_2=2xF_1.
\end{equation}
$q_f(x)\dd x$ 是适当因子化定义下的数分布, 不可当作所有尺度下固定的经典概率. QCD 辐射使分布随 $Q^2$ 演化, 产生标度破坏. 胶子虽电中性, 仍通过分裂和辐射影响结构函数和强子碰撞.

\originalsection{渐近自由与禁闭}
一圈耦合满足 $\dd\alpha_s/\dd\log Q^2=-\beta_0\alpha_s^2/(4\pi)$, $\beta_0=11-2n_f/3$ (SU(3)). 积分给
\begin{equation}
 \alpha_s(Q^2)=\frac{4\pi}{\beta_0\log(Q^2/\Lambda_{\rm QCD}^2)}.
\end{equation}
$\beta_0>0$ 时高尺度耦合变小, 这是渐近自由. 对 $Q$ 接近 $\Lambda$ 本式失去微扰控制, 不能把公式的极点当作精确物理奇点. 禁闭是低能非微扰性质, 不能仅凭一圈对数式证明.

喷注是高能 parton 级联与强子化后形成的集中能流. 质子对撞的总硬截面需把 parton 截面同两侧 PDF 卷积:
\begin{equation}
 \sigma_{pp\to X}=\sum_{ij}\int\dd x_1\dd x_2\,f_i(x_1,\mu)f_j(x_2,\mu)
 \widehat\sigma_{ij\to X}(x_1x_2s,\mu).
\end{equation}
质子束流能量不能全部用于某次 parton 碰撞, $\widehat s=x_1x_2s$. 因子化尺度、喷注定义和高阶修正都是比较预测与测量时的必要约定.
